\section{Finite lattice translations and the normalization of theta measures}
\label{lm-section}

We now pass from the one-dimensional ordinary theta laws to the
finite operations which occur in the lattice construction.  These
operations introduce two kinds of constants: square roots of
lattice indices and an algebraic factor correcting a determinant.
Both must be retained.  An absolute value of one over the complex
numbers does not imply an absolute value of one at a finite place.
We prove the required arithmetic normalization and all the finite
measure arguments.  The precise analytic identity which identifies
the submitted higher-dimensional theta function is stated separately.

\subsection{Integral moments in a lattice of directions}

Fix $p$, a complete field $L\supset\mathbb Q_p$, an abstract
embedding $\iota$ of the complex coefficients into $L$, and a
period $\omega\in L$.  For an entire function $F$ on a
finite-dimensional complex space $V$ and an additive subgroup
$\Lambda\subset V$, say that its period-scaled moments are
integral on $\Lambda$ if, for every finite family
$v_1,\ldots,v_r\in\Lambda$, the formal series
\[
 \widehat F\left(\omega\sum_{i=1}^r v_i\log(1+T_i)\right)
                 \quad\hbox{belongs to }\mathcal O_L[[T_1,\ldots,T_r]].
 \tag{LM1}
\]
The notation means transport of Taylor coefficients and of the
coordinates of the $v_i$ by the fixed embedding.  It makes no
claim about convergence in the original complex topology.
By Proposition~\ref{ot-logarithm}, (LM1) is equivalent to
existence of the associated measures of norm at most one.

\begin{lemma}\label{lm-directions}
Suppose the series associated to $v_1,\ldots,v_r$ has coefficient
norm at most $C$.  Let
$w_i=\sum_j a_{ji}v_j$ for $1\le i\le s$, where
$\iota(a_{ji})\in\mathbb Z_p$.  Then the series associated to
$w_1,\ldots,w_s$ has norm at most $C$.
Consequently integral moments are preserved by finite products,
finite linear combinations with integral transported coefficients,
and pullback by a complex linear map $g:V'\to V$ satisfying
\[
 \forall v\in\Lambda'\quad
 \exists N\in\mathbb Z_{>0},\ p\nmid N,\quad Ngv\in\Lambda.
 \tag{LM2}
\]
\end{lemma}
\begin{proof}
In the original series substitute
\[
 T_j=\prod_{i=1}^s(1+S_i)^{\iota(a_{ji})}-1.
 \tag{LM3}
\]
These are integral series with zero constant term.  Every
coefficient of the substituted series is a finite sum of original
coefficients times integral scalars, so its absolute value is at
most $C$.  The formal identity
$\log\prod_i(1+S_i)^{a_i}=\sum_i a_i\log(1+S_i)$ follows by
differentiation and the constant term, and identifies the series
with the stated directions.  This proves the first assertion.
Products and finite integral linear combinations follow by
multiplying or adding the series for one common direction family.
For a finite family in $\Lambda'$, choose the finitely many
integers in (LM2).  Apply the first assertion to the vectors
$Ngv\in\Lambda$ with diagonal coefficients $1/N\in\mathbb Z_p$.
The chain rule identifies this result with the pullback.
\end{proof}

The same proof permits coefficients in the ring of integers of a
split imaginary quadratic field: their images lie in $\mathbb Z_p$.
In particular every finite family of directions in
$\mathcal O_k$ can be reduced to the direction $1$ in the
one-dimensional theorem of Section~\ref{ct-section}.  Products of
the one-dimensional theta functions therefore have integral
moments on the product lattice, with the same period.

\subsection{The finite translation operator}

Let $H$ be a positive definite Hermitian form on $V$, linear in
its first variable, and put $E=\operatorname{Im}H$.  For $x\in V$
define
\[
 (U_xF)(z)=e^{-\pi H(z,x)-\pi H(x,x)/2}F(z+x).
 \tag{LM4}
\]
Then $U_xU_y=e^{-\pi iE(x,y)}U_{x+y}$.
Let $\Lambda_0,\Lambda_1$ be commensurable full lattices on which
$E$ takes integer values.  Suppose a function $\alpha$ on their
rational span restricts to a semicharacter on each lattice:
\[
 \alpha(x+y)=\alpha(x)\alpha(y)e^{\pi iE(x,y)},\qquad
 |\alpha(x)|=1,\qquad x,y\in\Lambda_j.
 \tag{LM5}
\]
The values on the intersection are the same because a single
function $\alpha$ is used.  A theta function for
$(H,\Lambda_j,\alpha)$ is an entire function satisfying
$U_lF=\alpha(l)^{-1}F$ for $l\in\Lambda_j$.
Put $W_x=\alpha(x)U_x$ and
$d=[\Lambda_1:\Lambda_0\cap\Lambda_1]$.  With $R$ any set of
representatives of this finite quotient, define
\[
 I_{\Lambda_0,\Lambda_1}F
                 =\frac1{\sqrt d}\sum_{x\in R}W_xF,
 \tag{LM6}
\]
where $\sqrt d$ is the positive real square root before transport.

\begin{proposition}\label{lm-finite-intertwiner}
The expression (LM6) is independent of the representatives and
is a theta function for $(H,\Lambda_1,\alpha)$.
Suppose moreover $p\nmid d$, every $\alpha(x)$ on the rational
lattice is a root of unity, and the period-scaled moments of
$U_tF$ on $\Lambda_0$ are integral whenever
$Nt\in\Lambda_0$ for some $N>0$ prime to $p$.
Then $I_{\Lambda_0,\Lambda_1}F$ and all its prime-to-$p$
translates have integral moments on $\Lambda_1$.
\end{proposition}
\begin{proof}
For $x,y\in\Lambda_1$, (LM4)--(LM5) and the integrality of
$E(x,y)$ give $W_xW_y=W_{x+y}$.  If
$l\in\Lambda_0\cap\Lambda_1$, then $W_lF=F$.
Thus replacing $x$ by $x+l$ does not change $W_xF$.
Translation by any $y\in\Lambda_1$ permutes the quotient, so
$W_yI_{\Lambda_0,\Lambda_1}F=I_{\Lambda_0,\Lambda_1}F$.
This proves the first assertion, including entireness.

The finite group of order $d$ is killed by $d$, so
$d\Lambda_1\subset\Lambda_0$.  Each representative $x$ has
prime-to-$p$ order modulo $\Lambda_0$.  If
$Nt\in\Lambda_1$ with $p\nmid N$, then
$dN(t+x)\in\Lambda_0$ for every $x\in R$.  Also
\[
 U_t(W_xF)=\alpha(x)e^{-\pi iE(t,x)}U_{t+x}F.
 \tag{LM7}
\]
The pairing $E$ is rational on the common rational span, hence
the exponential in (LM7) is a root of unity.
The transported coefficient $1/\sqrt d$ has norm one since its
square is $1/d$; every root of unity also has norm one.
The assumed translated integrality, (LM2) with $g=1$ and
$N=d$, and Lemma~\ref{lm-directions} now prove the assertion.
\end{proof}

Two successive operators of this form are therefore finite
linear combinations of translates with integral coefficients
whenever both indices are prime to $p$.  Expanding the two sums
introduces only the cocycle in (LM4).  Multiplication by an
algebraic character value which is a $p$-adic unit preserves
the conclusion.  These observations apply to the submitted
construction's two successive lattice moves; their validity
does not depend on choosing representatives with small
archimedean size.

\subsection{An arithmetic factor correcting the determinant}

Let $k$ be an imaginary quadratic field, with class number $h$,
and let $(W,h_W)$ be a Hermitian $k$-space isometric to the
standard Hermitian space $k^m$.  Fix a full
$\mathcal O_k$-lattice $M\subset W$ and one rational isometry
$\gamma_0:W\to k^m$.  For each rational split prime $\ell$ in
a set $\mathcal S$, assume that there is an isometry
$\gamma_\ell$ carrying
$M\otimes\mathbb Z_\ell$ onto
$\mathcal O_k^m\otimes\mathbb Z_\ell$.
Choose one prime $\mathfrak l\mid\ell$, a generator
\[
       \mathfrak l^h=(\pi_\ell),\qquad
       r_\ell^{\,2h}=\frac{\pi_\ell}{\bar\pi_\ell},
 \tag{LM8}
\]
in an algebraic closure, and put
\[
 y_\ell=v_{\mathfrak l}
             \bigl(\det(\gamma_\ell\gamma_0^{-1})\bigr).
 \tag{LM9}
\]
The valuation is normalized to have value one on a uniformizer.
The existence of $\pi_\ell$ follows directly from the definition
of the class number, and the root $r_\ell$ is nonzero.

\begin{theorem}\label{lm-kappa}
The integer $y_\ell$ is independent of the locally integral
isometry $\gamma_\ell$.  Only finitely many of these integers
are nonzero, so
\[
                         \kappa=\prod_{\ell\in\mathcal S}
                                         r_\ell^{-y_\ell}
 \tag{LM10}
\]
is a well-defined nonzero algebraic number.
If $p\in\mathcal S$ and $\gamma$ is any isometry satisfying
$\gamma(M\otimes\mathbb Z_p)=\mathcal O_k^m\otimes\mathbb Z_p$
(so both maps between these specified lattices are integral), then,
for every embedding $\iota:\overline{\mathbb Q}\to\mathbb C_p$,
\[
 |\iota(\kappa)|_p^2
  \left|\iota\bigl(\det(\gamma\gamma_0^{-1})\bigr)\right|_p=1.
 \tag{LM11}
\]
Moreover every complex conjugate of $\kappa$ has absolute value one.
\end{theorem}
\begin{proof}
If $\gamma$ and $\gamma'$ both carry the completed lattice
onto the standard lattice, then $\gamma\gamma'^{-1}$ and its
inverse are integral matrices over
$\mathcal O_k\otimes\mathbb Z_\ell$.  Its determinant is a
unit at both primes above $\ell$.  Taking valuations proves
the independence in (LM9).  There is one nonzero integer $B$
such that both $\gamma_0M$ and $\mathcal O_k^m$ are contained
in $B^{-1}$ times the other.  For $\ell\nmid B$ their
completions are equal, and $\gamma_0$ itself is an allowable
local isometry.  Therefore $y_\ell=0$ away from the finite
set of prime divisors of $B$.

If $\ell\ne p$, the ideals of $\pi_\ell$ and $\bar\pi_\ell$
are supported above $\ell$.  Their images at every place
above $p$ are units.  Equation (LM8) gives
$|\iota(r_\ell)|_p=1$, so (LM10) yields
\[
                   |\iota(\kappa)|_p
                        =|\iota(r_p)|_p^{-y_p}.
 \tag{LM12}
\]
Set $a=\det(\gamma\gamma_0^{-1})$.
The isometry equation implies $a\bar a=1$.
At the chosen prime $\mathfrak p$ above $p$,
$v_{\mathfrak p}(a)=y_p$ by the independence already proved.
At its conjugate the valuation is $-y_p$.
Since $p$ splits, both ramification indices are one, while
$v_{\mathfrak p}(\pi_p)=h$ and
$v_{\bar{\mathfrak p}}(\pi_p)=0$.
If $\iota$ induces $\mathfrak p$, (LM8) consequently gives
\[
 |\iota(a)|_p=p^{-y_p}
                   =|\iota(r_p)|_p^{2y_p}.
 \tag{LM13}
\]
If it induces $\bar{\mathfrak p}$, both exponents change sign
and the same equality holds.  Multiplying by (LM12) squared
proves (LM11).
Finally every complex embedding of $k$ is either the given
one or its conjugate.  Thus $\pi_\ell/\bar\pi_\ell$ has
complex absolute value one in every embedding; every conjugate
of its $2h$-th root has absolute value one as well.  The finite
product (LM10) has the stated archimedean property.
\end{proof}

\begin{corollary}\label{lm-change-scalar}
Suppose $a$ is the determinant in (LM11), and $c$ is an
algebraic number satisfying $c^8=a^4$.  Then
$|\iota(\kappa c)|_p=1$.
\end{corollary}
\begin{proof}
Taking absolute values in the eighth-power identity gives
$|\iota(c)|_p^2=|\iota(a)|_p$.
The assertion is now (LM11).
\end{proof}

This is the exact exponent needed in the source's change of
theta normalization.  Merely knowing that $c$ or $\kappa$
has complex absolute value one would not give this conclusion.
The proof also explains why the choice of one prime over each
split rational prime and the sign in $r_\ell^{-y_\ell}$ matter.

\subsection{The precise fourth tensor criterion}

Let $\Lambda_0=\mathcal O_k^m\subset\mathbb C^m$, with its
product Hermitian form and product theta function
$\Theta_0(z)=\prod_{j=1}^m\theta(z_j)$ as in
Section~\ref{ct-section}.  Let $V$ have complex dimension
$m/4$, and let $\Lambda\subset V$ be a full lattice stable
under $\mathcal O_k$.  Suppose
$g:V^4\to\mathbb C^m$ is a complex linear isometry, and
$M_g=g(\Lambda^4)$ is commensurable with $\Lambda_0$.
Assume a common semicharacter on their rational span, with
root-of-unity values, and write
$j_g=[M_g:M_g\cap\Lambda_0]$.  Define the entire function
\[
 K_g(z)=j_g^{-1/2}
       \sum_{x\in M_g/(M_g\cap\Lambda_0)}
                       \alpha(x)(U_x\Theta_0)(gz).
 \tag{LM14}
\]
As above the expression uses any finite set of representatives.

\begin{theorem}\label{lm-fourth-criterion}
Suppose $p\ge5$ is split in $k$, the one-dimensional theta
function has the ordinary laws of Section~\ref{ct-section},
and $F$ is an entire function on $V$.  Suppose there are two
isometries $g_0,g_1$ as above and algebraic nonzero numbers
$\kappa,c$ such that
\[
 \begin{aligned}
 F^{\otimes4}&=\kappa K_{g_0},&
 K_{g_0}&=cK_{g_1},&
 c^8&=\det(g_1g_0^{-1})^4,\\
 p&\nmid j_{g_1},&
 |\iota(\kappa)|_p^2
      |\iota(\det(g_1g_0^{-1}))|_p&=1 .
 \end{aligned}
 \tag{LM15}
\]
Then every prime-to-$p$ translate $U_tF$ has integral
period-scaled moments on $\Lambda$.
\end{theorem}
\begin{proof}
Corollary~\ref{lm-change-scalar}, or its one-line norm
calculation, gives $|\iota(\kappa c)|_p=1$.
The product function $\Theta_0$ and all its prime-to-$p$
translates have integral moments on $\Lambda_0$, by the
one-dimensional theorem and Lemma~\ref{lm-directions}.
Proposition~\ref{lm-finite-intertwiner} proves the same assertion
for $I_{\Lambda_0,M_{g_1}}\Theta_0$ on $M_{g_1}$.

If $Nt\in\Lambda$, with $p\nmid N$, put
$t^{[4]}=(t,t,t,t)$.  Since $g_1$ is an isometry,
direct substitution in (LM4) gives
\[
 U_{t^{[4]}}(K_{g_1})
       =(U_{g_1t^{[4]}}
              I_{\Lambda_0,M_{g_1}}\Theta_0)\circ g_1.
 \tag{LM16}
\]
The vector $g_1t^{[4]}$ has prime-to-$p$ order modulo
$M_{g_1}$.  Pullback by $g_1$, which carries the direction
lattice $\Lambda^4$ exactly onto $M_{g_1}$, preserves integral
moments.  Applying (LM15) and (LM16) gives integral moments
for
\[
 U_{t^{[4]}}F^{\otimes4}=(U_tF)^{\otimes4}.
 \tag{LM17}
\]
Theorem~\ref{ot-tensor-root} now proves the assertion for $U_tF$.
\end{proof}

The arithmetic part of the last condition of (LM15) is precisely
Theorem~\ref{lm-kappa} when $g_0,g_1$ arise by complexifying the
rational isometries there.  The remaining analytic statement
$K_{g_0}=cK_{g_1}$ with its eighth-power determinant relation
is a substantive theta transformation theorem.  It is not proved
by the elementary moment criterion, by linearity of derivatives,
or by the complex norm of the two functions.

\subsection{Finite Fourier projection and support of measures}

Let $X=\mathbb Z_p^r$, let $\mu$ be a bounded $L$-valued
measure, and suppose $L$ contains all $p$-power roots of unity.
For a continuous function $w$, the measure $w\mu$ means
$(w\mu)(f)=\mu(wf)$.  For $n\ge1$ and $a\in X$, choose a
primitive $p^n$-th root $\zeta_n\in L$.

\begin{lemma}\label{lm-character-projection}
There is an exact equality of continuous functions
\[
 {\bf1}_{a+p^nX}(x)=p^{-nr}
     \sum_{y\in(\mathbb Z/p^n\mathbb Z)^r}
                       \zeta_n^{\,\langle y,x-a\rangle}.
 \tag{LM18}
\]
Consequently, restriction of $\mu$ to any residue class has
norm at most $\|\mu\|$, even though the displayed expression
contains a denominator divisible by $p$.
If
\[
       \chi_y\mu=\chi_y(a)\mu
       \quad\hbox{for every character }\chi_y
               \hbox{ of }X/p^nX,
 \tag{LM19}
\]
then $\mu$ is supported on $a+p^nX$.
\end{lemma}
\begin{proof}
The sum factors as a product of $r$ geometric sums.
A one-dimensional sum is $p^n$ if $x_i-a_i$ is zero
modulo $p^n$, and zero otherwise: in the latter case multiply
the sum by the nontrivial root $\zeta_n^{x_i-a_i}$ and subtract.
This proves (LM18).  Its left side has supremum norm one,
which proves the norm assertion.
Multiplying $\mu$ by (LM18) and applying (LM19) gives $\mu$
again.  For any continuous $f$ vanishing on the residue class,
$\mu(f)=\mu({\bf1}_{a+p^nX}f)=0$.  This is the stated support.
\end{proof}

Thus estimating the terms of the Fourier sum separately would
introduce a spurious loss of $p^{nr}$.  The projection itself
has integral values and has norm one.  The same proof applies
to the inverse image of a coset in any finite quotient of $X$,
using the characters of that quotient.

\begin{theorem}\label{lm-finite-approximation}
Let $Y\subset X$ be closed and let $D\subset Y$ be dense.
Suppose $\mu$ is supported on $Y$ and $\|\mu\|\le C$.
For every finite family $f_1,\ldots,f_s\in C(X,L)$ and
$\epsilon>0$, there are $d_1,\ldots,d_t\in D$ and
$a_1,\ldots,a_t\in L$ with $|a_i|\le C$ such that
\[
 \left|\mu(f_j)-\sum_{i=1}^t a_if_j(d_i)\right|<\epsilon
                             \qquad(1\le j\le s).
 \tag{LM20}
\]
If $C=1$, all coefficients are integral.
\end{theorem}
\begin{proof}
For $C=0$ the measure is zero and the empty sum suffices.
For $C>0$, uniform continuity gives $n$ for which every $f_j$
oscillates by less than $\epsilon/C$ on each residue class
modulo $p^n$.  There are finitely many such classes meeting
$Y$.  Each has a point $d_i\in D$, since its intersection
with $Y$ is relatively open and nonempty.  Let $B_i$ be the
class and set $a_i=\mu({\bf1}_{B_i})$.
Then $|a_i|\le\|\mu\|\le C$.
The function equal to $f_j(d_i)$ on each of these classes
differs from $f_j$ by less than $\epsilon/C$ on their union.
Outside that union, the measure vanishes by its support
condition.  Applying its norm bound proves (LM20).
The case $Y=\varnothing$ again gives the zero measure.
\end{proof}

This proves the finite simultaneous approximation required by a
polynomial-kernel argument as soon as the actual translation
identities establish the support and the rational lattice is
dense in that support.  It preserves the norm bound on each
coefficient and does not insert a factor equal to the number
of points in the finite sum.

\subsection{Limits of the reduction and counterexamples}

The preceding results prove every finite operation in the
passage from a locally normalized fourth tensor identity to
integral lattice moments, and the exact finite Fourier
projection used to obtain support.  To identify them with the
submitted global kernels one still has to supply the theta
transformation in (LM15), the locally standard isometries with
the required period lattices, and the actual higher-dimensional
translation equations which impose the stated kernel support.
Their names in the formal source do not discharge these
mathematical obligations.

Two false shortcuts are worth excluding.  First, over
$k=\mathbb Q(i)$ the number $a=(3+4i)/5$ has complex
absolute value one.  At the embedding into $\mathbb Q_5$
with $i\equiv2\pmod5$, its numerator is a unit, so
$|a|_5=5$.  The determinant correction in (LM11) cannot be
replaced by an archimedean norm calculation.
Second, multiplication of a theta function by $p^{-1}$
preserves its transformation law but destroys an integral
constant term of norm one.  The theta transformation law
alone therefore does not fix the arithmetic normalization.
These examples refute weakened claims, not the statements
with the normalizations retained above.


